\exnumber{03}
\extitle{UI III - Ex 4 - Accessibilité (application)}
\exfiles{ex03/fichier Figma, nomme Name\_FirstName\_DeliverableName\_Date\_VersionNumber}
\exauthorize{Those named in the exercise}
\makeheaderfiles

\textbf{Instructions :}

Choose an app you find particularly accessible. Find at least 3 concrete examples of accessible measures and explain how they ease the reading or perception of visual content.

\textbf{Livrables :}

\begin{itemize}
\item A 3-page Figma file with 3 different elements from an app (it can be from 3 different apps) that are accessible. The three measures must be different.
\item Déposez le fichier Figma sur Github, nommé ainsi : « Name\_FirstName\_DeliverableName\_Date\_VersionNumber ». Exemple : « Doe\_John\_MidFiPrototype\_05242024\_V1 ».
\item Ajoutez un titre dans le fichier.
\item Soignez l'aspect global du fichier. Il doit être clair, simple et lisible. Vous pouvez regarder canva.com, sans surcharger le dessin de détails.
\end{itemize}

\textbf{Ressources :}

\begin{itemize}
\item What accessibility is and why it's so important \url{https://uxdesign.cc/what-accessibility-is-and-why-its-so-important-9c56e033ff26}
\item Web Content Accessibility Guidelines (WCAG) \url{https://www.w3.org/WAI/standards-guidelines/wcag/}
\item Accessibility testing  ---  W3C Wik \url{https://www.w3.org/wiki/Accessibility_testing}i
\item UX Myth: Accessibility is expensive and difficult \url{https://uxmyths.com/post/654091803/myth-5-accessibility-is-expensive-and-difficult}
\item Salesforce UX - 7 tips every designer needs to know about accessibility \url{https://medium.com/salesforce-ux/7-things-every-designer-needs-to-know-about-accessibility-64f105f0881b}
\end{itemize}

\textbf{Critères :}

\begin{itemize}
\item Examples of measures
\begin{itemize}
\item Color contrast
\item Highlighted or enlarged text when hovering
\item Links color that is different from the body text's
\item Breadcrumbs to indicate where you are in terms of navigation
\item In forms, the title of the case should be above the text zone
\item Error states
\item etc.
\end{itemize}
\end{itemize}

